% Generated draft: requires mathematical and source review.
\documentclass[11pt]{article}
\usepackage[margin=25mm]{geometry}
\usepackage{fontspec}
\setmainfont{Latin Modern Roman}
\usepackage{amsmath,amssymb,mathtools,unicode-math}
\setmathfont{Latin Modern Math}
\usepackage{enumitem,needspace,xurl,hyperref}
\setlength{\parindent}{0pt}
\setlength{\parskip}{5pt}
\setlength{\emergencystretch}{3em}
\newcommand{\N}{\mathbb{N}}
\newcommand{\Z}{\mathbb{Z}}
\newcommand{\Q}{\mathbb{Q}}
\newcommand{\R}{\mathbb{R}}
\newcommand{\C}{\mathbb{C}}
\newcommand{\F}{\mathbb{F}}
\newcommand{\A}{\mathbb{A}}
\newcommand{\PP}{\mathbb{P}}
\newcommand{\E}{\mathbb{E}}
\newcommand{\eps}{\varepsilon}
\newcommand{\dd}{\,\mathrm{d}}
\begin{document}
\section*{UnsolvedMath ID: 30006117}
\textbf{Source problem code:} catalogue entry
\subsection{Definitions and mathematical statement}
Let \(G\) be a quasi-split group and let \((U,\psi)\) be a Whittaker datum, in the setting intended by the cited report. The supplied extract does not specify the base field, the class of groups under consideration, the coefficient ring \(\Lambda\), or the precise definitions of the geometric objects and categories below.

Write \(\operatorname{Par}_G\) for the parameter object denoted by that symbol in the source, and \(\operatorname{Isoc}_G\) for its object of isocrystals. Their precise constructions are not supplied. The notation \(\operatorname{Perf}(\operatorname{Par}_G)\) denotes the category of perfect complexes on the intended parameter object. The notation \(\operatorname{Coh}_{\operatorname{Nilp}}(\operatorname{Par}_G)\) is understood as the intended category of coherent objects with nilpotent singular support; the relevant support theory and nilpotent locus require specification from the source. Likewise, \(\operatorname{Shv}(\operatorname{Isoc}_G,\Lambda)\) denotes the intended category of \(\Lambda\)-valued sheaves, and the superscript \(\omega\) is understood to select its compact objects. These interpretations of the categorical notation require confirmation against the full report.

\textbf{Conjecture 2 (as supplied).} For every quasi-split \(G\) and Whittaker datum \((U,\psi)\) satisfying the intended source hypotheses, there is a canonical \(\operatorname{Perf}(\operatorname{Par}_G)\)-linear equivalence
\[
\operatorname{Coh}_{\operatorname{Nilp}}(\operatorname{Par}_G)
\simeq
\operatorname{Shv}(\operatorname{Isoc}_G,\Lambda)^{\omega}.
\]
Here linearity means compatibility with the intended actions of \(\operatorname{Perf}(\operatorname{Par}_G)\) on the two categories. Those actions, the categorical framework, and the normalization or dependence on \((U,\psi)\) that makes the equivalence canonical are not specified in the extract. Consequently, this records the source's conjectural equivalence with its unresolved prerequisites; it is not a fully specified standalone formulation.
\subsection{Short English statement}
The conjecture asks for a canonical equivalence between the spectral category with nilpotent support and compact objects in the sheaf category on isocrystals, compatible with the action of perfect complexes on the parameter object. The base field, coefficients, categorical constructions, and Whittaker normalization remain unspecified in the supplied extract.
\subsection{Sources}
\begin{itemize}
\item Local UnsolvedMath registry, supplied record with permanent ID \(30006117\), exact source code \texttt{catalogue entry}, entitled \emph{Spectral Categorical Langlands for Quasisplit Groups}.
\item Supplied report citation: \emph{Representations of \(p\)-adic Groups} (2025), Oberwolfach Reports, DOI 10.4171/owr/2024/54, pp.\ 3103--3178. \url{https://doi.org/10.4171/owr/2024/54}. The extract identifies the statement as Conjecture 2 in \emph{Spectral action on isocrystals}, attributed to Arnaud Eteve, joint work with Dennis Gaitsgory, Alain Genestier, and Vincent Lafforgue. The registry separately records proposed year 2024.
\item Supplied literature source, OWR report: \url{https://publications.mfo.de/bitstream/handle/mfo/4270/OWR_2024_54.pdf?sequence=1}.
\item Supplied literature source, labeled ``Zhu, tame categorical LLC'': \url{https://arxiv.org/abs/2504.07482}.
\item Supplied literature source, labeled ``Schematic and analytic LLC'': \url{https://arxiv.org/abs/2606.02799}.
\item Literature assessment dated \texttt{2026-08-22}: the supplied curation labels the problem partially solved and states that the full spectral categorical equivalence remains conjectural, while nontrivial tame and unipotent pieces have been constructed or proved and current work targets broader cases. This is supplied dated curation, not an independently verified assessment; its time-sensitive status and cited claims require independent verification.
\item The registry labels the statement \texttt{corrected\_verified} and reports that it was reconstructed and checked against the official Oberwolfach report. This verification claim is supplied metadata and has not been independently checked for this draft.
\end{itemize}
\end{document}
